% CLAUDE REPORT::WCD-PVEF-0001::2026-11-13T16:52:00-08:00
% !TEX program = xelatex
% Compiles with XeLaTeX or pdfLaTeX (ASCII source).
\documentclass[11pt]{article}
\usepackage[margin=1in]{geometry}
\usepackage{amsmath,amssymb,amsthm}
\usepackage{xurl}
\newtheorem{theorem}{Theorem}
\newtheorem{lemma}[theorem]{Lemma}
\newtheorem{proposition}[theorem]{Proposition}
\newtheorem{corollary}[theorem]{Corollary}
\theoremstyle{remark}
\newtheorem*{remark}{Remark}
\newcommand{\SAT}{\mathrm{SAT}}
\newcommand{\EF}{\mathrm{EF}}
\newcommand{\ER}{\mathrm{ER}}
\newcommand{\PV}{\mathrm{PV}}
\newcommand{\Wb}{\overline{W}}
\newcommand{\poly}{\mathrm{poly}}
\newcommand{\RejSound}{\mathrm{RejSound}}
\newcommand{\status}[1]{\textbf{[#1]}}

\title{CLAUDE WCD-PVEF-0001\\[2pt]
\large The worst-case derivation bridge $B3$ for Extended-Frege-strength systems:\\
Cook's PV translation, the one-sided certification hypothesis, and the exact residual gap}
\author{CLAUDE scheduled routine (worst-case minimum-derivation directive)}
\date{2026-11-13 16:52 America/Los\_Angeles}

\begin{document}
\maketitle

\noindent\textbf{Status.} \textbf{P versus NP remains OPEN.} This report proves no lower bound and
no unconditional bridge. It anchors the corpus's primary conditional bridge $B3$ to a classical
known theorem for Extended-Frege-strength systems, isolates the exact residual hypothesis, and maps
current primary literature onto it. Additive report in the CLAUDE lane. It changes no schedules,
task count, TMR-1..5 role ownership, root \texttt{main.tex}, \texttt{proof1.txt} or \texttt{proof2.txt}.

\section{Target (verbatim from the governing directive)}
Fix one sound, complete, polynomial-time-checkable refutation system $R$. Then
\[
 W_R(F)=\min_{\pi:F\vdash_R\bot}\mathrm{cost}_R(\pi),\qquad
 \Wb_R(N)=\max_{\substack{F\in\mathrm{UNSAT}\\ |\langle F\rangle|\le N}} W_R(F).
\]
The inner minimum grants free selection of the best derivation, so it lower-bounds the cost of
\emph{executing} any $R$-derivation, independently of how the derivation is discovered.
The primary bridge (2026-11-02 directive, $B3$) is
\[
 \SAT\in\mathrm P\ \Longrightarrow\ \exists C,k\ \forall N:\ \Wb_R(N)\le C(N+1)^k
 \qquad\text{for the same preselected }R. \tag{B3($R$)}
\]
$\Wb_R$ superpolynomial together with B3($R$) excludes $\SAT\in\mathrm P$.

\section{Corpus reviewed}
Read this run: \texttt{AGENTS.md}; the TMR replacement directive
\texttt{research/coordination/2026-11-13\_1249\_TM\_RESOURCE\_AGENT\_REPLACEMENT.md};
the 2026-11-02 worst-case minimum-derivation refocus
\texttt{research/coordination/2026-11-02\_1204\_WORST\_CASE\_MINIMUM\_DERIVATION\_FIVE\_ROLE\_REFOCUS.md}
(B1/B2/B3 definitions); \texttt{proof1.txt} and \texttt{proof2.txt} heads and the B3-related
entries [P2-B2], [P2-L5-409], [P2-L5-424]; the Codex audit
\texttt{research/agents/codex/2026-11-11\_1203\_SAT\_POLYNOMIAL\_ASSUMPTION\_AUDIT.tex};
LSPV-1/337 (finite reflection bridge), LSPV-1/415 ($\Pi^0_1$-completeness of clocked rejection
soundness), LSPV-1/592 (conditional ER/EF simulation strength of the native calculus) and its
handoff; MAIN/REFEREE331 and /343 (vacuity of B3 when $\SAT\notin\mathrm P$); the TMR-1..5
role directories and handoffs (latest commits 2026-11-13 16:29 PST). Repository search for
\texttt{Krajicek}, \texttt{PV}, \texttt{Pich}, \texttt{Santhanam}, \texttt{reflection principle}.

\section{Duplication / contradiction audit}
\begin{itemize}
\item \emph{Already in corpus, adopted:} B3 holds iff $R$ size-simulates the identity system $Q_A$ of a
 hypothetical polynomial decider $A$ (LSPV-5/409); B3 is vacuously true if $\SAT\notin\mathrm P$, and
 for any $R$ with an \emph{unconditionally known} superpolynomial $\Wb_R$ (ordinary Resolution:
 Haken 1985, Urquhart 1987, Chv\'atal--Szemer\'edi 1988) B3($R$) is \emph{equivalent} to
 $\SAT\notin\mathrm P$ (REFEREE331/343); finite-length reflection proofs in a theory $T$ suffice
 (LSPV-1/337); no r.e.\ theory certifies \emph{every} correct clocked decider (LSPV-1/415); a native
 calculus with unrestricted paid AND2 definitions and cross-parent Resolution p-simulates ER, hence its
 lower-bound half is at least as hard as an ER/EF lower bound (LSPV-1/592, conditional on A1--A4).
\item \emph{Not found in corpus (new here):} the explicit anchoring of B3 for EF-strength $R$ to
 Cook's 1975 PV-to-EF translation; the observation that only \emph{one-sided} (rejection) soundness
 needs to be PV-provable while completeness need only be true; the resulting exact residual
 hypothesis PV-CERT; and the map from Pich--Santhanam 2023 to sufficient conditions for B3(EF).
\item \emph{No contradiction} with any verified entry. In particular the result below does not
 conflict with LSPV-1/415: that theorem forbids certifying \emph{all} correct programs, while
 PV-CERT asks for \emph{one}.
\end{itemize}

\section{Results}

Let $\PV$ be Cook's equational theory of polynomial-time functions. For a polynomial-time
(explicitly clocked) program $A$ with $\PV$ function symbol $f_A$, and the $\PV$ symbol
$\mathrm{Sat}(F,a)$ for CNF evaluation, define the open (universal) $\PV$ formula
\[
 \RejSound_A(F,a):\qquad f_A(F)=\mathtt{UNSAT}\ \rightarrow\ \mathrm{Sat}(F,a)=0 .
\]
Say $A$ is \emph{UNSAT-complete} if $A(F)=\mathtt{UNSAT}$ for every well-formed unsatisfiable $F$.

\begin{theorem}[Cook 1975 translation applied to B3; \status{KNOWN THEOREM, assembly}]\label{thm:pvef}
Let $A$ be polynomial-time and UNSAT-complete. If $\PV\vdash\RejSound_A$, then there is a polynomial
$p$ (depending on $A$ and the fixed $\PV$ proof) such that every unsatisfiable CNF $F$ with
$|\langle F\rangle|\le N$ has an Extended Resolution refutation of size at most $p(N)$.
Hence $\Wb_{\ER}(N)\le p(N)$ and, by Cook--Reckhow p-equivalence, $\Wb_{\EF}(N)\le\poly(N)$.
\end{theorem}
\begin{proof}
Cook's translation \cite{cook75} gives, for each input length $n$, EF proofs $\pi_n$ of the
propositional translations $\|\RejSound_A\|_n$ of size $\poly(n)$, constructible in polynomial time,
over atoms for the bits of $F$, the bits of $a$, and extension atoms naming the polynomial-size
circuits that compute $f_A$ and $\mathrm{Sat}$. Fix an unsatisfiable $F$ of length $\le N$.
(i) Substituting the constant bits of $F$ for the $F$-atoms in every line of $\pi_n$ yields an EF
proof of the restricted formula of the same size up to a constant factor (EF is closed under
substitution of constants for atoms, extension definitions being instantiated accordingly).
(ii) Because $A$ is UNSAT-complete, the circuit for $f_A$ on the constant input evaluates to
$\mathtt{UNSAT}$; EF derives, gate by gate, the value of every extension atom of that variable-free
circuit, hence $\|f_A(F)=\mathtt{UNSAT}\|$, in size polynomial in the circuit size.
(iii) Modus ponens gives $\|\mathrm{Sat}(F,a)=0\|$ with free atoms $a$ only. With $F$ fixed, EF
proves in polynomial size the equivalence of $\|\mathrm{Sat}(F,a)=1\|$ with $\bigwedge_j C_j(a)$, so
from hypotheses $C_1(a),\dots,C_m(a)$ it derives $\bot$. (iv) Cook--Reckhow \cite{cookreckhow79}:
EF and ER p-simulate each other on CNF refutations. All overheads are polynomial in $N$, with the
polynomial fixed by $A$ and its $\PV$ proof.
\end{proof}

\begin{remark}[one-sidedness; \status{PROVED/EXACT}]
Only rejection soundness enters the $\PV$ proof. UNSAT-completeness is used solely in step (ii), as a
\emph{true} fact about the concrete computation, which EF verifies by evaluation. Neither acceptance
soundness nor completeness needs to be provable, and acceptance behaviour on satisfiable inputs is
irrelevant. This is weaker than ``$\PV$ proves $\mathrm{P}=\mathrm{NP}$'' in the sense of
Cook--Kraj\'{\i}\v{c}ek \cite{ck07}, which requires the full correctness statement.
\end{remark}

Define the hypothesis
\[
 \text{PV-CERT:}\quad \SAT\in\mathrm P\ \Longrightarrow\ \exists A\ \text{polytime, UNSAT-complete, with }\PV\vdash\RejSound_A .
\]

\begin{corollary}[\status{PROVED/EXACT} from Theorem~\ref{thm:pvef} and LSPV-1/592]\label{cor:b3}
(a) PV-CERT implies B3(EF) and B3(ER).
(b) If a fixed native calculus $R$ satisfies the LSPV-1/592 hypotheses A1--A4 (unrestricted paid AND2
definitions with exact facets, cross-parent Resolution, faithful weights; currently CONDITIONAL / not
authenticated for historical $R_{\mathrm{mix341}}$), then $\Wb_R\le\poly(\Wb_{\ER})$, so PV-CERT
implies B3($R$).
(c) Contrapositive: if $\Wb_{\EF}$ is superpolynomial, then no polynomial-time UNSAT-complete $A$ has
a $\PV$ proof of $\RejSound_A$. This is the classical ``consistency of $\mathrm{P}\neq\mathrm{NP}$ with
$\PV$'' consequence (stated for full correctness in \cite{ck07}), not $\mathrm P\neq\mathrm{NP}$.
\end{corollary}

\begin{proposition}[theory ladder; \status{PROVED/EXACT} schema, \status{OPEN} instances]
For every consistent r.e.\ theory $T\supseteq\PV$ with polynomial-time proof checking there is a
sound complete Cook--Reckhow system $R_T$ (EF plus the axiom scheme of $T$-provable reflection,
as in LSPV-1/337 and \cite{ps23} \S1.1.2) such that
$T\text{-CERT}\Rightarrow$ B3($R_T$). Raising $T$ weakens the certification hypothesis but strengthens
$R_T$, so the required lower bound on $\Wb_{R_T}$ becomes at least as hard. No fixed $T$ is known to
satisfy $T$-CERT; LSPV-1/415 shows no single $T$ certifies every correct clocked decider.
\end{proposition}
\begin{proof}
Same argument as Theorem~\ref{thm:pvef}, with the reflection axiom for the $T$-proof of
$\RejSound_A$ replacing Cook's translation; the proof of $\RejSound_A$ is a fixed constant-size object
and its propositional instances are verified in $R_T$ by substitution and evaluation.
\end{proof}

\section{Exact residual gap}
For EF-strength $R$ the worst-case derivation route now reads
\[
 \underbrace{\Wb_{\EF}\ \text{superpolynomial}}_{\text{OPEN since 1979}}
 \ \wedge\ \underbrace{\text{PV-CERT}}_{\text{OPEN}}
 \ \Longrightarrow\ \SAT\notin\mathrm P .
\]
PV-CERT fails exactly in a world where $\mathrm P=\mathrm{NP}$ but every polynomial-time SAT decider
has rejection soundness independent of $\PV$. Excluding that world is the single missing bridge
arrow for this route. For weak $R$ (Resolution, cutting planes, bounded-depth Frege, polynomial
calculus) the lower bound is known and the bridge is \emph{equivalent} to the conclusion
(REFEREE331/343), so those systems provide no reduction of difficulty. Free selection of the best
derivation does not change either statement: it is already built into $W_R$.

\section{Literature check (this run)}
\begin{enumerate}
\item Cook (1975), PV and its EF translation \cite{cook75}: \status{KNOWN THEOREM}; used in
 Theorem~\ref{thm:pvef}. Directly transferable.
\item Cook--Reckhow (1979) \cite{cookreckhow79}: a p-bounded system exists iff NP=coNP; EF and ER
 p-equivalent. \status{KNOWN THEOREM}. Existential over systems; does not certify a chosen $R$.
\item Cook--Kraj\'{\i}\v{c}ek (2007) \cite{ck07}: $\PV$ proves P=NP iff $\PV$ proves NP=coNP iff $\PV$ defines a
 polynomial-time algorithm returning a satisfying assignment or an EF refutation; superpolynomial EF
 lower bounds would make $\mathrm P\ne\mathrm{NP}$ consistent with $\PV$. \status{KNOWN THEOREM}.
\item Pich--Santhanam (FOCS 2019) \cite{ps19}: under hardness assumptions there are natural
 distributions of hard tautologies for which proving lower bounds for strong enough systems is
 infeasible. \status{KNOWN THEOREM, conditional}. Barrier for the lower-bound half.
\item Pich--Santhanam (arXiv 2312.08163, 2023) \cite{ps23}: if EF is not p-bounded, then
 $\mathrm P\ne\mathrm{NP}$ follows from any of: $S^1_2$ proves that a p-time function generates
 anticheckers for SAT (formal Theorem~7); $S^1_2$ proves that a p-time function witnesses an error of
 each p-size circuit failing on SAT; or an OWF/hardness-of-E condition (formal Theorem~8); results
 extend to theories $T\supseteq S^1_2$ with systems $P_T$. Abstract and introduction checked this
 run; full proofs not re-verified. \status{KNOWN THEOREM, conditional}. These are the \emph{only}
 literature-supplied sufficient conditions found for B3(EF) other than PV-CERT, and they are
 provability hypotheses of the same type.
\item J.~Li, ECCC TR25-187 (2025) \cite{li25}: a time-complexity hierarchy for PV proofs measured by
 EF size of translations; a proposed conjecture implying NP=coNP unprovable in PV. Useful analogy for
 quantifying the polynomial in Theorem~\ref{thm:pvef}; no transfer theorem to B3 yet
 (ANALOGY/CANDIDATE).
\end{enumerate}
Cross-domain transfer check: the structure ``soundness proof in a weak theory $\Rightarrow$ short
certificates'' is the proof-carrying-code / certified-compilation pattern (CompCert-style translation
validation). Classification: useful structural analogy; the exact transfer theorem is Cook's
translation, already used above. Post-quantum ledger: special-soundness/extraction (item 6) matches
the shape of step (ii) only as analogy; no cryptographic assumption is used.

\section{Hostile controls}
\begin{itemize}
\item \emph{Tseitin / Gaussian elimination:} polynomial-time decidable, $2^{\Omega(n)}$ Resolution
 (Urquhart). Falsifies pointwise B1 for Resolution; consistent with Theorem~\ref{thm:pvef} because
 Gaussian elimination has PV-provable rejection soundness and Tseitin has polynomial EF refutations.
\item \emph{Pigeonhole:} exponential for Resolution (Haken), polynomial for EF (Cook--Reckhow);
 consistent.
\item \emph{Clocked padding / sentinel controls (LSPV-1/337, /415, /462):} Theorem~\ref{thm:pvef} uses
 one fixed $A$ and one fixed PV proof; no per-length proof, oracle or advice; the polynomial is fixed
 before $N$. Quantifiers checked: $\exists A\,\exists p\,\forall N\,\forall F$.
\item \emph{Vacuity:} if $\SAT\notin\mathrm P$, PV-CERT and B3 hold vacuously; Corollary~\ref{cor:b3}(c)
 is stated as a contrapositive about provability only, so no vacuous strengthening is claimed.
\end{itemize}

\section{What changed relative to root main.tex}
Root \texttt{main.tex} not edited (TMR-5 owns whole-document synthesis under the 2026-11-13 directive;
this lane owns no region). Proposed addition for TMR-1/TMR-5 review: record PV-CERT as the named
residual arrow of the B3 route for EF-strength $R$, beside [P2-B2], with Theorem~\ref{thm:pvef} as
KNOWN THEOREM assembly.

\section{Adopted / rejected / retired components}
Adopted: B1/B2/B3 hierarchy; LSPV-5/409 identity simulation; LSPV-1/337, /415, /592; REFEREE331/343.
Rejected: treating B3 for Resolution-strength systems as a weaker sufficient bridge (it is equivalent
to the conclusion). Retired: none.

\section{Stale or overstated claims found}
None found in the entries reviewed. Caution added: ``P=NP provable in PV'' and PV-CERT are
hypotheses about provability; neither follows from $\SAT\in\mathrm P$ alone.

\section{Status table}
\begin{center}
\begin{tabular}{ll}
Theorem~\ref{thm:pvef} (PV-certified decider $\Rightarrow$ poly $\Wb_{\EF}$) & KNOWN THEOREM (assembly)\\
One-sided soundness remark & PROVED/EXACT\\
Corollary~\ref{cor:b3}(a),(c) & PROVED/EXACT\\
Corollary~\ref{cor:b3}(b) for historical $R_{\mathrm{mix341}}$ & CONDITIONAL on A1--A4\\
PV-CERT & OPEN\\
$\Wb_{\EF}$ superpolynomial & OPEN\\
B3 for any fixed $R$ & OPEN\\
P versus NP & OPEN
\end{tabular}
\end{center}

\section{One next target}
Test PV-CERT against the strongest available evidence in either direction: decide whether the
Pich--Santhanam anticheckers hypothesis (an $S^1_2$-provable p-time antichecker generator) implies
PV-CERT or $S^1_2$-CERT directly, or exhibit why the two conditions are incomparable. Owner suggestion:
TMR-1 (dependency chain), with TMR-5 adversarial review.

\begin{thebibliography}{9}
\bibitem{cook75} S.~A.~Cook. Feasibly constructive proofs and the propositional calculus. STOC 1975, 83--97.
\bibitem{cookreckhow79} S.~A.~Cook, R.~A.~Reckhow. The relative efficiency of propositional proof systems. J.~Symbolic Logic 44(1), 1979, 36--50.
\bibitem{ck07} S.~A.~Cook, J.~Kraj\'{\i}\v{c}ek. Consequences of the provability of NP $\subseteq$ P/poly. J.~Symbolic Logic 72(4), 2007, 1353--1371. \url{https://doi.org/10.2178/jsl/1203350791}
\bibitem{ps19} J.~Pich, R.~Santhanam. Why are proof complexity lower bounds hard? FOCS 2019, 1305--1324. \url{https://doi.org/10.1109/FOCS.2019.00080}
\bibitem{ps23} J.~Pich, R.~Santhanam. Towards P$\neq$NP from Extended Frege lower bounds. arXiv:2312.08163, 2023. \url{https://arxiv.org/abs/2312.08163}
\bibitem{li25} J.~Li. On the time complexity of feasible proofs. ECCC TR25-187, 2025. \url{https://eccc.weizmann.ac.il/report/2025/187/}
\end{thebibliography}

\end{document}
% END_OF_VERIFIED_REPORT::CLAUDE::WCD-PVEF-0001
